\chapter{Conclusion and Outlook}
\label{ch:chapter07}

\section{This is possible}
\label{sec:chapter06:possible}
Tools like BurpSuite, ZAP and SQLMap already implement partially automated penetration testing techniques through signature-driven \ac{DAST} scanning approaches. \ac{AI} models trained through Reinforcement Learning are used in these tools.

With the rise of \ac{LLM}s into the public more and more research was conducted on leveraging \ac{LLM}s for automated penetration testing. Research revealed that \ac{LLM}s are already capable of performing penetration tests based on their inherit knowledge. Without any technical additions publicly available \ac{LLM}s can conduct penetration tests. By providing an \ac{LLM} with tool usage capabilities the process can be automated nearly entirely. Additional domain specific knowledge provided to the \ac{LLM} improves the process further.

Corporate \ac{LLM} driven penetration testing tools like XBOW are able to find vulnerabilities in corporate tech stacks.

\section{Functional limitations}
\label{sec:chapter06:limitations}
XBOW's high ranking on the Hackerone leaderboard clearly proves that generally \ac{LLM} powered tools are capable of finding vulnerabilities in real world bug bounty programs. \cite{hackerone_leaderboard} \cite{xbow_hype} When it comes to performing penetration tests in proprietary tech stacks or applications built with niche frameworks, automated penetration tests leveraging \ac{LLM}s reach their limits. Tests with proprietary tools like XBOW revealed that even highly sophisticated tools struggle with niche and proprietary tech stacks. 


\section{Future Work}
\label{sec:chapter06:future_work}
With an appropriate setup, \ac{LLM}s can conduct end-to-end penetration tests. Their vulnerability identification performance remains limited on proprietary or niche technology stacks. Future research should prioritize investigating the underlying causes of these limitations and developing strategies to enhance \ac{AI}-driven testing in corporate environments. Additionally, further further research should be conducted to address and mitigate the security challenges associated with these approaches.